\chapter{Maintenance, Troubleshooting, and Extension}
\label{ch:maintenance}

\chapteroverview
  {Defines the complete change cycle, package impact rules, extension recipes,
   diagnosis workflow, data policy, and rollback discipline.}
  {Start with the standard change cycle, select the matching extension recipe,
   diagnose with read-only probes before mutation, and finish with verified
   documentation, evidence, and rollback.}

\section{Standard Change Cycle}

\begin{procedurebox}[One complete maintenance loop]
\begin{enumerate}
  \item Record branch, commit, worktree state, environment, and the command or
        observation that demonstrates the issue.
  \item Classify the owner: interface, description/plugin, control/config, or
        bringup policy. Classify the data as product source, local output,
        external, or immutable.
  \item Find all consumers with \code{rg}; include filenames, parameter keys,
        message fields, topic names, PDDL symbols, hashes, and left/right aliases.
  \item Make the smallest coherent source change. Never repair \file{install/}
        or mutate a promoted/frozen artifact.
  \item Run syntax/schema and focused source tests, then inspect the diff.
  \item Build the affected package dependency closure and source the new overlay.
  \item Verify the installed surface, registered tests, and relevant headless/
        GUI/runtime scenario.
  \item Update public documentation, this manual, and the Repository Source
        Index.
  \item Record exact evidence, exclusions, rollback, and any migration action.
\end{enumerate}
\end{procedurebox}

Before editing, treat all existing worktree changes as maintainer-owned. Do not
discard or reformat unrelated work. Use version-control review to separate the
intended patch from local build/capture output.

\section{Change Impact by Package}

\begin{longtable}{L{0.23\textwidth} L{0.30\textwidth} L{0.38\textwidth}}
\caption{Package-level change impact and compatibility risk}\label{tab:package-change-impact}\\
\toprule
\tableheader Owner changed & Rebuild/review downstream & Typical compatibility risk \\
\midrule
\endfirsthead
\toprule
\tableheader Owner changed & Rebuild/review downstream & Typical compatibility risk \\
\midrule
\endhead
\file{mfja_rail_interfaces} & Both control and bringup dependency closure & Generated bindings, serializers, recordings, field semantics \\ \rowseparator
\file{mfja_3rd_floor_description} & Description; control/bringup when path/plugin contract changes & SDF ABI, world service name, joint/topic alignment, asset URI/license \\ \rowseparator
\file{mfja_robot_control_config} & Control plus bringup when wrapper-visible & Install list, parameter precedence, runtime contracts, research artifacts \\ \rowseparator
\file{mfja_3rd_floor_bringup} & Bringup & Public arguments, profile defaults, singleton/startup order \\
\bottomrule
\end{longtable}

\noindent\textbf{Governing paths:} each package's root
\file{package.xml} and \file{CMakeLists.txt}, plus dependency declarations in
the downstream package manifests. Exact package roots are listed in
Table~\ref{tab:executive-package-map}.

This matrix identifies downstream ownership and compatibility risk; it does
not replace the canonical change-to-test matrix in
Table~\ref{tab:change-test-matrix}.

Shared Python foundations deserve broader regression. The dependency hot-spot
table in the Repository Source Index (\cref{app:catalog}) identifies modules
with many in-repository importers. In particular, treat rail defaults/fleet,
visual dataset/contracts, runtime contracts, PDDL scenario generation, and task
execution as cross-cutting.

\section{Extension Recipes}

Table~\ref{tab:extension-source-map} is the source registry for every recipe in
this section. It keeps paths in one place so the procedural text does not
repeat them inconsistently.

\begin{longtable}{L{0.22\textwidth} L{0.68\textwidth}}
\caption{Extension recipe source-path registry}\label{tab:extension-source-map}\\
\toprule
\tableheader Change type & Authoritative source paths \\
\midrule
\endfirsthead
\toprule
\tableheader Change type & Authoritative source paths \\
\midrule
\endhead
Python tool/node & \file{mfja_robot_control_config/scripts/}; \file{mfja_robot_control_config/CMakeLists.txt}; \file{mfja_robot_control_config/package.xml}; \file{mfja_robot_control_config/test/} \\ \rowseparator
Launch surface & \file{mfja_3rd_floor_bringup/launch/}; \file{mfja_robot_control_config/launch/}; \file{mfja_3rd_floor_bringup/CMakeLists.txt}; \file{mfja_robot_control_config/CMakeLists.txt} \\ \rowseparator
Installed configuration & \file{mfja_robot_control_config/config/}; \file{mfja_robot_control_config/CMakeLists.txt} \\ \rowseparator
ROS interface & \file{mfja_rail_interfaces/msg/}; \file{mfja_rail_interfaces/srv/}; \file{mfja_rail_interfaces/CMakeLists.txt}; consumer code under \file{mfja_robot_control_config/} \\ \rowseparator
Gazebo model/world/robot & \file{mfja_3rd_floor_description/models/}; \file{mfja_3rd_floor_description/worlds/}; \file{mfja_3rd_floor_description/urdf/}; \file{mfja_3rd_floor_description/src/}; \file{mfja_3rd_floor_description/THIRD_PARTY.md} \\ \rowseparator
Rail geometry/routing/devices & \file{mfja_robot_control_config/config/room_315_kinematics/}; \file{mfja_robot_control_config/scripts/room_315_kinematic_shuttle.py}; \file{mfja_robot_control_config/scripts/room_315_kinematic_shuttle_node.py} \\ \rowseparator
TaskGoal, PDDL, and execution & \file{mfja_robot_control_config/config/room_315_vla/pddl/}; \file{mfja_robot_control_config/scripts/room_315_task_goal_*.py}; \file{mfja_robot_control_config/scripts/room_315_pddl_*.py}; \file{mfja_robot_control_config/scripts/room_315_task_execution*.py} \\ \rowseparator
Dataset, model, and deployment & \file{mfja_robot_control_config/config/room_315_visual_state/}; \file{mfja_robot_control_config/config/room_315_vla/}; relevant \file{mfja_robot_control_config/scripts/room_315_vla_*.py}; \file{mfja_robot_control_config/scripts/room_315_build_runtime_candidate_v4.py}; \file{mfja_robot_control_config/scripts/room_315_promote_runtime_v4.py}; \file{mfja_robot_control_config/scripts/room_315_authorize_runtime_v4_closed_loop.py} \\
\bottomrule
\end{longtable}

\subsection{Add a Python tool or ROS node}

\begin{enumerate}
  \item Put the file under \file{mfja_robot_control_config/scripts}, with a
        Python 3 shebang, concise module docstring, validated inputs, meaningful
        exit status, and safe \code{--help}/\code{--dry-run} where it mutates.
  \item Separate ROS-free logic from the node adapter when practical.
  \item Add it to the explicit CMake \code{install(PROGRAMS)} list only when it
        should be an installed/public program; set the source executable bit.
  \item Declare required package/system dependencies; optional research imports
        must fail with an actionable prerequisite message.
  \item Add focused tests and register suitable deterministic tests in CMake.
  \item Rebuild and verify \code{ros2 pkg executables}, copied install, help,
        failure cases, and output/side effects.
\end{enumerate}

Do not flatten another large multi-purpose module into the already broad script
directory without a clear family/contract. Its Repository Source Index category
and public CLI status must remain understandable.

\subsection{Add a launch file or argument}

Define units, legal choices, default, and side effects. Trace it through every
wrapper/include; test default, explicit override, false/zero/empty values, and
parameter type. Use \code{--show-args} after rebuild. If the argument deletes a
file, downloads data, starts actuation, or alters an authorization boundary,
state that before the operator command.

When both the VLA supervisor and visual runtime are composed, assign exactly one
camera bridge. When overriding VLA command/status topics, change gateway topics
to match. Keep obstacle flags semantically consistent across floor and gateway.

\subsection{Add an installed config}

Adding a file to Git does not guarantee CMake installs it. Add the appropriate
install rule, validate with its owning loader, rebuild, then inspect
\code{ros2 pkg prefix --share mfja_robot_control_config}. Test a normal copied
install as well as \code{--symlink-install}. Do not install a whole directory
when ignored/local files could leak into the artifact.

\subsection{Add or change a ROS interface}

Define unit, authority, legal/empty/stale semantics, and compatibility policy.
Update every consumer and recorded representation. Rebuild the interface first,
then dependent packages; verify generated definitions with
\code{ros2 interface show}. String-valued modes/actions should be centralized
and validated; for a future breaking version, consider message constants or
enums rather than another unrestricted string.

\subsection{Add a Gazebo model, world asset, or robot}

Preserve upstream license/provenance and exact license text. Validate model
URI, mesh scale/axis, pose, material, rendering cost, and collision intent.
Add \file{model.config}; include the model from the appropriate world; update
the Repository Source Index. For a robot, align SDF/URDF joint names, bridge
topics, plugin limits, registry aliases, TF prefix, state feedback, and commands.

Keep visual-only assets explicitly collision-free. Do not silently add heavy
mesh collisions to industrial robots; that changes performance and behavior and
needs a deliberate collision/control validation plan.

\subsection{Change rail geometry, routing, or devices}

Use the mapping in \cref{ch:configuration}. Preserve CSV originals and migrate
explicitly; compare length, endpoint, tangent, route cycles, guard failure, and
all device ratios. Test both sides and left public/internal conversion. Verify
Gazebo world alignment with markers and representative shuttles, including
headway, stopper crossing, target clamp, occupied reset, and pose-service fault.

\subsection{Change task semantics or PDDL}

Treat this as one dependency chain:

\begin{lstlisting}
TaskGoal schema/validation -> grounding -> problem generator
  -> expert and runtime domains -> PlanSys2 response
  -> full-plan translator/validator -> supervisor primitive
  -> fresh effect verification -> tests/documentation
\end{lstlisting}

Unsupported compound meaning must still clarify or reject. Add invalid and
adversarial cases, not only a new happy path. Preserve first-primitive execution
and replan-after-effect behavior.

\subsection{Change datasets, models, or deployment policy}

Create new versioned directories/manifests. Never overwrite frozen roles,
sealed final tests, checkpoint sidecars, promotion manifests, authorizations,
or evidence. Run split/leakage audit, record seeds/tool commit, qualify on CPU
and target device, shadow, acceptance, and fault campaigns, then perform manual
promotion and separate task authorization. Retain and test the complete prior
bundle for rollback.

\section{Troubleshooting by Symptom}

Use read-only observation first. The source groups below are ordered by startup
dependency: floor, robots, rails, camera/visual state, planning, then
installation and retained data.

\begin{longtable}{L{0.10\textwidth} L{0.23\textwidth} L{0.57\textwidth}}
\caption{Troubleshooting source-group registry}\label{tab:troubleshooting-sources}\\
\toprule
\tableheader Key & Subsystem & Authoritative source paths \\
\midrule
\endfirsthead
\toprule
\tableheader Key & Subsystem & Authoritative source paths \\
\midrule
\endhead
D1 & Floor launch/world services & \file{mfja_3rd_floor_bringup/launch/room_315_floor_common.py}; \file{mfja_robot_control_config/launch/multi_robot_sim.launch.py}; \file{mfja_robot_control_config/launch/room_315_launch_utils.py} \\ \rowseparator
D2 & Robots and plugins & \file{mfja_robot_control_config/config/robots.yaml}; \file{mfja_robot_control_config/launch/multi_robot_sim.launch.py}; \file{mfja_3rd_floor_description/src/} \\ \rowseparator
D3 & Rail runtime & \file{mfja_robot_control_config/launch/room_315_dual_kinematic_shuttles.launch.py}; \file{mfja_robot_control_config/scripts/room_315_kinematic_shuttle_node.py}; \file{mfja_robot_control_config/config/room_315_kinematics/} \\ \rowseparator
D4 & Camera and visual state & \file{mfja_robot_control_config/launch/room_315_vla_supervisor.launch.py}; \file{mfja_robot_control_config/scripts/room_315_visual_runtime_v4.py}; \file{mfja_robot_control_config/config/room_315_vla/visual_state_runtime.yaml} \\ \rowseparator
D5 & Planning and task gateway & \file{mfja_robot_control_config/launch/room_315_task_execution.launch.py}; \file{mfja_robot_control_config/scripts/room_315_task_execution_node.py}; \file{mfja_robot_control_config/config/room_315_vla/task_execution_runtime.yaml} \\ \rowseparator
D6 & Install, tests, and retained data & \file{mfja_3rd_floor_bringup/CMakeLists.txt}; \file{mfja_3rd_floor_description/CMakeLists.txt}; \file{mfja_rail_interfaces/CMakeLists.txt}; \file{mfja_robot_control_config/CMakeLists.txt}; \file{mfja_robot_control_config/test/}; \file{mfja_robot_control_config/scripts/room_315_vla_dataset_recorder.py}; \file{docs/project_documentation/Makefile} \\
\bottomrule
\end{longtable}

\begin{longtable}{L{0.20\textwidth} L{0.28\textwidth} L{0.34\textwidth} L{0.08\textwidth}}
\caption{Troubleshooting symptoms, checks, actions, and source groups}\label{tab:troubleshooting}\\
\toprule
\tableheader Symptom & Check & Likely cause/action & Source \\
\midrule
\endfirsthead
\toprule
\tableheader Symptom & Check & Likely cause/action & Source \\
\midrule
\endhead
Second floor launch refuses & Owning process and lock; ROS node list & Supported singleton is active. Stop it cleanly; do not delete/bypass a live lock & D1 \\ \rowseparator
No \rosname{/clock} & Gazebo server output, bridge process, sourced domain & Server/bridge failed, wrong overlay/domain, or delayed startup & D1 \\ \rowseparator
Clock present but motion frozen & Clock stamp changing; Gazebo pause state & Full floor starts paused; unpause or relaunch running & D1 \\ \rowseparator
World services absent & Internal world name and \rosname{/world/...} list & Filename/entity mismatch, bridge delay/failure, wrong world & D1 \\ \rowseparator
Robot absent & Registry selector, create response, Gazebo entity list & Unknown/disabled selector, bad model URI, spawn failure & D2 \\ \rowseparator
Joint command has no effect & Exact namespace/type/subscriber, simulation time, limits & Wrong selector/overlay, plugin not loaded, paused world, invalid target & D2 \\ \rowseparator
Rail topic absent & Four-second delay, side enable/count, rail-node logs & High-level rail disabled, node validation failed, stale overlay & D3 \\ \rowseparator
Shuttle is WAITING & stopper actual state, headway, route guard & Expected guard, not an OFF acknowledgement; clear safely or command OFF & D3 \\ \rowseparator
Shuttle is FALLING & current segment and actual switches; network successor & Missing/mismatched guarded transition; stop, clear occupancy, repair/reset deliberately & D3 \\ \rowseparator
Logical state moves, visual does not & Set-pose service/client diagnostics and entity name & Service timeout/failure or missing entity; preserve logs, reconcile/restart & D3 \\ \rowseparator
Visual orphan after REMOVE & Gazebo entity and delete response & Logical removal preceded failed asynchronous delete; restart/explicit expert cleanup & D3 \\ \rowseparator
No camera image & bridge count, Gazebo sensors/rendering/EGL, topic type & Duplicate/missing bridge, renderer failure, wrong world/topic & D4 \\ \rowseparator
Visual runtime rejects startup & Manifest path/hash/schema/deployment mode/environment & Missing/incomplete/wrong external bundle; do not weaken validation & D4 \\ \rowseparator
Observations stay unready & Both images, both rail heartbeats, freshness, supervisor status, diagnostics & Paused clock, missing known-empty heartbeat, stale sync, safety gate & D4 \\ \rowseparator
Planner inactive & \code{ros2 lifecycle get /planner}, launch logs & Delayed configure/activation, missing PlanSys2 dependency, disabled bundled planner & D5 \\ \rowseparator
Gateway cannot plan & GetPlan service, runtime PDDL, fresh observation/status & Planner absent, invalid problem/goal, unready authorized inputs & D5 \\ \rowseparator
Execution remains disabled & Config, exact authorization/promotion/checkpoint hashes & Intended fail-closed behavior; restore/qualify exact bundle & D5 \\ \rowseparator
Executable/config missing after build & CMake install list, source mode, package share, overlay prefix & Source-only file, non-executable symlink, copied-install path bug, stale overlay & D6 \\ \rowseparator
Test count unexpectedly small & CMake cache/Torch import, registered list, pytest exit code & Optional test registration or source-only tests; reconfigure and record scope & D6 \\ \rowseparator
Disk usage grows & Workspace/repo \file{log/}, datasets, external bundles & Preserve needed evidence, then apply project retention policy; never delete unknown user data & D6 \\
\bottomrule
\end{longtable}

Useful read-only probes include:

\begin{lstlisting}[language=bash]
ros2 node list | sort
ros2 topic list -t | sort
ros2 service list -t | sort
ros2 topic info --verbose <topic>
ros2 param dump <node>
ros2 lifecycle get /planner
ros2 doctor --report
gz service --list
\end{lstlisting}

Compare types and QoS at both ends before adding sleeps or changing reliability.
Use timestamps to distinguish stale feedback from an incorrect value.

\section{Build, Temporary, and User Data}

\begin{longtable}{L{0.23\textwidth} L{0.34\textwidth} L{0.34\textwidth}}
\caption{Data classes, locations, and maintenance policy}\label{tab:data-lifecycle}\\
\toprule
\tableheader Class & Examples & Maintenance action \\
\midrule
\endfirsthead
\toprule
\tableheader Class & Examples & Maintenance action \\
\midrule
\endhead
Product source & package code/config/assets/docs & Review, test, version \\ \rowseparator
Build output & \file{build/}, \file{install/}, \file{log/}, Python/pytest caches & Recreate; never hand-edit or treat as only copy \\ \rowseparator
Runtime temporary & bridge YAML, patched SDF, filtered world under \file{/tmp} & Inspect during diagnosis; process-owned; may remain after shutdown \\ \rowseparator
User captures & \file{~/.ros/...datasets} and reports & Back up/retain by experiment policy; do not delete implicitly \\ \rowseparator
External runtime & checkpoints, GGUF, promotion/authorization bundles & Verify hashes/provenance; immutable after promotion \\ \rowseparator
Frozen research & split/final-test/evidence manifests & Supersede with new version, never rewrite \\ \rowseparator
Third-party asset & CAD/meshes/model derivatives & Preserve exact terms, attribution, originals \\
\bottomrule
\end{longtable}

\noindent\textbf{Classification sources:} package source roots listed in
Table~\ref{tab:executive-package-map}, temporary-path generation in
\file{mfja_robot_control_config/launch/multi_robot_sim.launch.py}, dataset
output handling in
\file{mfja_robot_control_config/scripts/room_315_vla_dataset_recorder.py}, and
third-party provenance in
\file{mfja_3rd_floor_description/THIRD_PARTY.md}.

Large local logs currently exist both inside this checkout and at workspace
level; neither is a source tree. Establish a documented retention policy before
cleanup, because failed-run logs and research evidence may be needed.

\section{Documentation Maintenance}

Keep detailed truth in one chapter and cross-reference it. Update this manual
when any package, executable, launch argument/default, public topic/service,
message field, config schema, authority rule, artifact gate, operator command,
or safety boundary changes.

From \file{docs/project_documentation/}:

\begin{lstlisting}[language=bash]
make check
\end{lstlisting}

The Repository Source Index covers Git-tracked source, so untracked experiments
do not silently become documented product. Inspect the PDF table of contents,
List of Tables, row separators, command wrapping, cross-references, and source
classification. The documentation build and validation rules are defined in
\file{docs/project_documentation/Makefile}.

\section{Release and Rollback}

For code/config rollback, apply a reviewed version-control revert to source,
rebuild the same dependency closure, source the intended overlay, and rerun the
failure-reproducing test plus smoke checks. Never roll back by editing an
installed file.

For learned runtime rollback, select the previous complete immutable promoted
bundle and its matching task authorization, update the host-local config,
verify all expected hashes, run shadow/readiness, and only then re-enable. A
checkpoint alone is not a rollback unit.

Preserve failing logs, manifests, input identity, commands, and results until
the cause and corrective action are reviewed.
